\Needspace{11\baselineskip}
\begin{center}
\textit{The number of ambiguous classes.}
\end{center}

\Article{257}
Of all classes in a given order of a given determinant, the ambiguous classes in particular call for a fuller investigation, and determining their number will open the way to many other things. It is enough, however, to determine this number in the properly primitive order alone, since the remaining cases can easily be reduced to this one. We will do this by first showing how to find all properly primitive ambiguous forms \((A,B,C)\) of the proposed determinant \(D\) in which either \(B=0\) or \(B=\frac12 A\), and then how to find from their number the number of all properly primitive ambiguous classes of determinant \(D\).

\SourceLeaf{290}
I. Clearly all properly primitive forms \((A,0,C)\) of determinant \(D\) are found by taking for \(A\) each divisor of \(D\) (both positive and negative) for which \(C=-\frac{D}{A}\) is relatively prime to \(A\). Thus when \(D=1\), there are two such forms, \((1,0,-1)\), \((-1,0,1)\); just as many when \(D=-1\), namely \((1,0,1)\), \((-1,0,-1)\); when \(D\) is a prime number or a power of a prime (with either a positive or a negative sign), there are four: \((1,0,-D)\), \((-1,0,D)\), \((D,0,-1)\), \((-D,0,1)\). In general, when \(D\) is divisible by \(n\) distinct primes (among which \(2\) too must be counted here), there are altogether \(2^{n+1}\) such forms. Namely, setting \(D=\pm PQR\ldots\), where \(P,Q,R\) etc. denote distinct primes or powers of distinct primes, whose number is \(=n\), the values of \(A\) are \(1,P,Q,R\) etc. and products of any number of these numbers. By the theory of combinations the number of these values is \(2^n\), but it must be doubled, since each value must be given both a positive and. a negative sign.

II. Likewise, all properly primitive forms \((2B,B,C)\) of determinant \(D\) are obtained by taking for \(B\) all divisors of \(D\) (positive and negative) for which \(C=\frac12(B-\frac{D}{B})\) is an integer relatively prime to \(2B\). Since \(C\) must be odd, and therefore \(CC\equiv1\ (\mathrm{mod.}\,8)\), it follows from \(D=BB-2BC=(B-C)^2-CC\) that \(D\) is either \(\equiv3\ (\mathrm{mod.}\,4)\) when \(B\) is odd, or \(\equiv0\ (\mathrm{mod.}\,8)\) when \(B\) is even. Thus whenever \(D\) is congruent to one of the numbers \(1,2,4,5,6\) modulo \(8\), there are no such forms. When \(D\equiv3\ (\mathrm{mod.}\,4)\), \(C\) is an odd integer whatever divisor of \(D\) is taken for \(B\); but for \(C\) to have no common divisor with \(2B\), \(B\) must be chosen so that \(\frac{D}{B}\) is relatively prime to \(B\). Thus for \(D=-1\) there are two forms, \((2,1,1)\), \((-2.\ -1,-1)\), and in general it is easy to see that if the number of all primes dividing \(D\) is \(n\), altogether \(2^{n+1}\) forms result. --- When \(D\) is divisible by \(8\), \(C\) is an integer if we take for \(B\) any even divisor of \(\frac12 D\); the other condition, that \(C=\frac12 B-\frac{D}{2B}\) be relatively prime to \(2B\), is satisfied \textit{first} by taking for \(B\) all divisors of \(D\) that are twice an odd number and for which \(\frac{D}{B}\) has no common divisor with \(B\). Their number (counting the different signs) is \(2^{n+1}\) if \(D\) is assumed divisible by \(n\) distinct odd primes; \textit{second}, by taking for \(B\) all divisors of
\SourceLeaf{291}
\(\frac12 D\) that are divisible by four and for which \(\frac{D}{2B}\) is relatively prime to \(B\). Their number too is \(2^{n+1}\), so in this case there are altogether \(2^{n+2}\) such forms. Namely, setting \(D=\pm2^\mu PQR\ldots\), where \(\mu\) denotes an exponent greater than \(2\) and \(P,Q,R\) distinct odd primes or powers of such primes whose number is \(n\), we can take \textit{both} for \(\frac12 B\) \textit{and} for \(\frac{D}{2B}\) the values \(1,P,Q,R\) etc. and products of any number of these numbers, with both positive and negative signs.

From all this we conclude that if \(D\) is assumed divisible by \(n\) distinct odd primes (setting \(n=0\) when \(D=\pm1\) or \(\pm2\) or a power of two), the number of all properly primitive forms \((A,B,C)\) in which \(B\) is either \(0\) or \(\frac12 A\) is \(2^{n+1}\) when \(D\) is either \(\equiv1\) or \(\equiv5\ (\mathrm{mod.}\,8)\); \(2^{n+2}\) when \(D\equiv2,3,4,6\text{ or }7\ (\mathrm{mod.}\,8)\); and finally \(2^{n+3}\) when \(D\equiv0\ (\mathrm{mod.}\,8)\). Comparing this with what we gave in Article 231 for the number of all possible characters of primitive forms of determinant \(D\), we see that the former is in every case exactly twice the latter. Moreover, when \(D\) is negative, there are clearly as many positive as negative forms among those forms.
